\section{Regras e Inferência}
\label{sec:regras}

A aplicação de regras é o coração do funcionamento do eProfile. Uma regra é uma instrução de inferência registrada por uma aplicação que será executada sobre um conjunto de trilhas de uma entidade, resultando num perfil de entidade de acordo com um modelo de usuário previamente cadastrado pela aplicação.

Para que uma regra possa ser aplicada, as seguintes informações precisam ter sido previamente registradas no eProfile:

\begin{itemize}

\item A regra propriamente dita;

\item Um modelo de usuário, que servirá como base para o formato de perfil da entidade;

\item A trilha da entidade, sobre a qual a inferência será executada.

\end{itemize}

A partir destas três informações, o raciocinador tem condições de gerar um perfil de usuário, de acordo com a Figura~\ref{fig:dados_inferencia}.

\begin{figure}
  \includegraphics[width=0.8\textwidth]{imagens/dados_inferencia.png}
  \caption{Fluxo de dados para inferência do perfil de usuário}
  \label{fig:dados_inferencia}
\end{figure}

As regras são sempre definidas pela aplicação, e uma regra é composto pelo seguinte conjunto de informações:

\begin{description}

\item[Nome:] um identificador único para a regra, que permitirá que a aplicação identifique a regra posteriormente, para que possa editá-la ou excluí-la;

\item[Campo do perfil:] define qual é o campo do perfil da entidade que deverá ser atualizado através da execução desta inferência. Este campo deve ser obrigatoriamente equivalente a uma classe definida na ontologia de modelo de usuário;

\item[Regra:] a regra que será executada sobre a trilha da entidade;

\item[Atualização:] uma instrução que define como o campo será atualizado (por exemplo, o valor pode ser simplesmente sobrescrito ou atualizado de acordo com alguma fórmula);

\item[Freqüência de atualização:] define de quanto em quanto tempo esta regra deve ser atualizada. É importante notar que o valor definido aqui pode não ser honrado, devido a outros motivos (como uma sobrecarga no servidor);

\item[Validade]: define qual a data limite para a execução desta regra. A ausência desta informação define que a regra continuará sendo processada indefinidamente;

\end{description}

O raciocinador gerenciará a carga de processamento no servidor. Assim, algumas regras podem não ser executadas de acordo com o agendamento solicitado. Além disso, é possível que uma aplicação registre certas regras mas deixe de usar o servidor, fazendo com que as regras continuem sendo executadas desnecessariamente. Assim, caso um perfil não seja consultado por mais de um mês, a freqüencia da execução das regras passa a diminuir progressivamente, até o ponto em que a regra não seja mais executada. Caso a aplicação volte a consultar os perfis, a freqüencia de execução das regras é normalizada.

As regras podem ser escritas em duas linguagens diferentes: SPARQL ou SWRL. O SPARQL \cite{sparql} é uma linguagem de consulta de ontologias cuja sintaxe lembra vagamente a linguagem de consulta de banco de dados SQL \cite{sql}, enquanto o SWRL \cite{swrl} é uma linguagem de regras que permite a consulta utilizando construtores lógicos mais complexos que o SPARQL.

Uma vez que a regra tenha sido validada pelo agente regulador, ela entra em funcionamento e passa a ser elegível para execução.

Tanto as consultas via OWL quanto as consultas via SPARQL são executadas através de um motor de inferências externo. Foram estudados os seguintes raciocinadores que oferecem suporte a estas operações: Pellet \cite{sirin05}, FaCT++ \cite{tsarkov06} e HermiT \cite{motik09}.

\begin{description}

\item[Pellet:] é um racionador com suporte a OWL DL, que oferece bom desempenho e um \emph{middleware} extensível. O Pellet é escrito em Java pode ser utilizado para raciocinar com tipos de dados definidos pelo usuário e provê suporte para depuração em ontologias \cite{sirin05}.

\item[FaCT++:] é um raciocinador de lógica proposicional criado como uma plataforma para experimentar com novos algoritmos de \emph{tableaux} semântico e técnicas de otimização. Ele incorpora a maior parte das técnicas de otimização padrão, mas inclui algumas novas \cite{tsarkov06}.

\item[HermiT:] é o primeiro raciocinador OWL que se baseia em um cálculo \emph{hypertableaux}, que provê desempenho superior a qualquer algoritmo anterior \cite{motik09}. O HermiT é disponibilizado como um plugin do Protegé, como uma ferramenta de linha de comando ou como uma biblioteca Java.

\end{description}

Devido ao conjunto de funcionalidades e ao desempenho superior, o HermiT foi escolhido como o raciocinador utilizado pelo eProfile.

\todo[inline]{Colocar um exemplo de uma regra}
